% !TEX root = ../matan.tex

\section{Критерий Коши для равномерной сходимости последовательностей. Пространства $\ell^\infty$, $C(K)$. Признак Вейерштрасса.}

\begin{Theorem}[required]{Критерий Коши для равномерной сходимости}{}
	Пусть $f_n \colon \Omega \to \C$. Последовательность $\{f_n\}$ равномерно сходится на $\Omega$ (к некоторой функции $f$) тогда и только тогда, когда она равномерно фундаментальна:
	\[
	\forall \eps > 0\ \exists N\ \forall n, m > N \colon \quad \sup_{x \in \Omega} |f_n(x) - f_m(x)| < \eps.
	\]
\end{Theorem}

\begin{proof}
	\textbf{Необходимость.} Пусть $f_n \to f$ равномерно. Для $\eps > 0$ найдём $N$ такое, что $\sup_x |f_n(x) - f(x)| < \eps/2$ при $n > N$. Тогда при $n, m > N$ для всех $x$
	\[
	|f_n(x) - f_m(x)| \le |f_n(x) - f(x)| + |f(x) - f_m(x)| < \tfrac{\eps}{2} + \tfrac{\eps}{2} = \eps,
	\]
	откуда $\sup_x |f_n(x) - f_m(x)| \le \eps$. % 

	\textbf{Достаточность.} Пусть выполнено условие Коши. Для каждого фиксированного $x$ числовая последовательность $\{f_n(x)\}$ фундаментальна (так как $|f_n(x) - f_m(x)| \le \sup_y|f_n(y)-f_m(y)|$), а $\C$ полно, поэтому существует предел $f(x) := \lim_n f_n(x)$; так задаётся функция $f \colon \Omega \to \C$.

	Покажем равномерность. Зафиксируем $\eps > 0$ и соответствующий $N$, так что $|f_n(x) - f_m(x)| < \eps$ для всех $x$ при $n, m > N$. Зафиксируем $n > N$ и $x$ и устремим $m \to \infty$; так как $f_m(x) \to f(x)$, в пределе получаем
	\[
	|f_n(x) - f(x)| \le \eps \qquad \forall x \in \Omega,\ \forall n > N.
	\]
	Значит, $\sup_x |f_n(x) - f(x)| \le \eps$ при $n > N$, то есть $f_n \to f$ равномерно.
\end{proof}

\begin{Definition}{Пространства $C(K)$ и $\ell^\infty$}{}
	\textbf{$C(K)$:} пусть $K$~--- компакт (например, $K \subset \RR^n$). Через $C(K)$ обозначают пространство непрерывных функций $f \colon K \to \mathbb{K}$ ($\mathbb{K} = \RR$ или $\C$) с равномерной нормой
	\[
	\|f\|_{C(K)} = \|f\|_{\infty} = \sup_{x \in K} |f(x)|.
	\]

	\textbf{$\ell^\infty$:} пространство ограниченных последовательностей $x = (x_1, x_2, \ldots)$, $x_k \in \mathbb{K}$, с нормой
	\[
	\|x\|_{\ell^\infty} = \sup_{k \in \NN} |x_k| < \infty.
	\]
	Оба пространства \textbf{банаховы}.
\end{Definition}

\begin{Remark}{Связь с критерием Коши}{}
	Норма $\|\cdot\|_\infty$ в обоих пространствах~--- это «супремум модуля», поэтому сходимость по норме в $C(K)$ совпадает с равномерной сходимостью функций, а в $\ell^\infty$~--- с равномерной по индексу $k$. Доказанный критерий Коши есть в точности утверждение о \emph{полноте} этих пространств: равномерно фундаментальная последовательность имеет (равномерный) предел. В $C(K)$ предел вдобавок непрерывен, так как равномерный предел непрерывных функций непрерывен, поэтому $C(K)$ замкнуто в $\ell^\infty(K)$ и тем самым полно.
\end{Remark}

\begin{Theorem}{Признак Вейерштрасса}{}
	Пусть $f_k \colon \Omega \to \C$. Если существуют числа $a_k \ge 0$ такие, что
	\[
	|f_k(x)| \le a_k \quad \forall x \in \Omega, \qquad \text{и ряд } \sum_{k=1}^{\infty} a_k \text{ сходится,}
	\]
	то функциональный ряд $\sum_{k=1}^{\infty} f_k$ сходится равномерно (и абсолютно) на $\Omega$.
\end{Theorem}

\begin{proof}
	Обозначим частичные суммы $S_n = \sum_{k=1}^{n} f_k$. Так как числовой ряд $\sum a_k$ сходится, по критерию Коши для числовых рядов для любого $\eps > 0$ найдётся $N$ такое, что при $m > n > N$ выполнено $\sum_{k=n+1}^{m} a_k < \eps$. Тогда для всех $x \in \Omega$
	\[
	|S_m(x) - S_n(x)| = \Bigl| \sum_{k=n+1}^{m} f_k(x) \Bigr| \le \sum_{k=n+1}^{m} |f_k(x)| \le \sum_{k=n+1}^{m} a_k < \eps.
	\]
	Следовательно, $\sup_{x \in \Omega} |S_m(x) - S_n(x)| \le \eps$ при $m > n > N$, то есть последовательность $\{S_n\}$ равномерно фундаментальна. По критерию Коши для равномерной сходимости она сходится равномерно, а значит, ряд $\sum f_k$ сходится равномерно на $\Omega$.
\end{proof}
